\section{Tiên đề tập lũy thừa}

\

Tiên đề tập lũy thừa phát biểu rằng với mọi tập hợp $x$, tồn tại một tập hợp $p$ mà nó chứa tất cả và chỉ chứa các phần tử là các tập hợp con của $x$. Viết dưới dạng kí hiệu:
\begin{equation*}
    \defMath{\forall x, \exists p, \forall y, \left( y \in p \iff y \subseteq x \right)}.
\end{equation*}
Tập $p$ thỏa mãn điều này được gọi là \defText{tập lũy thừa} của $x$, kí hiệu là $\defMath{\mathfrak{P}(x)}$, $\defMath{\mathcal{P}(x)}$, hay $\defMath{2^x}$\footnote{Cần phải để ý, các kí hiệu ở đây là một sự lạm dụng. Kí hiệu $\mathfrak{P}(x)$ không phải để biểu diễn vị từ, và $2^x$ không biểu diễn phép lũy thừa.}.

Cho trước tập hợp $x$, giả sử tồn tại hai tập lũy thừa của $x$ là $\mathfrak{P}(x)$ và $\mathfrak{B}(x)$. Chúng ta sẽ chứng minh rằng $\mathfrak{P}(x) = \mathfrak{B}(x)$. Giống như các chứng minh trước, chúng ta xuất phát từ định nghĩa để đưa ra:
\begin{equation*}
    \begin{cases}
        \forall y, \left( y \in \mathfrak{P}(x) \iff y \subseteq x \right) \\
        \forall y, \left( y \in \mathfrak{B}(x) \iff y \subseteq x \right)
    \end{cases}.
\end{equation*}
Từ đó, có được ngay, $\forall y, (y \in \mathfrak{P}(x) \iff y \in \mathfrak{B}(x))$, suy ra $\mathfrak{P}(x) = \mathfrak{B}(x)$.

\exercise Cho $X$ và $Y$ là hai tập hợp. Chứng minh rằng,
\begin{multicols}{2}
    \begin{enumerate}
        \item $\varnothing \in \mathfrak{P}(X)$;
        \item $\mathfrak{P}(\{X; Y\}) = \{\varnothing; \{X\}; \{Y\}; \{X; Y\}\}$;
        \item $X \subseteq Y \implies \mathfrak{P}(X) \subseteq \mathfrak{P}(Y)$;
        \item $\mathfrak{P}(X \cup Y) \supseteq \mathfrak{P}(X) \cup \mathfrak{P}(Y)$.
    \end{enumerate}
\end{multicols}

\solution

\setcounter{subexercise}{1}
\arabic{subexercise}. Có tập rỗng là tập con của mọi tập hợp. Cho nên $\varnothing \subseteq X$. Nhắc lại định nghĩa của tập lũy thừa, $\mathfrak{P}(X)$ là tập chứa tất cả các tập con của $X$. Do đó, $\varnothing \in \mathfrak{P}(X)$. Điều phải chứng minh.

Một hệ quả tự nhiên của tính chất này đó là tập lũy thừa thì không rỗng. Biểu diễn dưới dạng kí hiệu:
$$
    \forall w, \mathfrak{P}(w) \neq \varnothing.
$$

\stepcounter{subexercise}
\arabic{subexercise}. Tập rỗng là tập con của mọi tập hợp, nên $\varnothing \subseteq \{X; Y\}$, dẫn đến $\varnothing \in \mathfrak{P}(\{X; Y\})$.

Theo cách xây dựng của đôi không thứ tự, $X \in \{X; Y\}$. Tập hợp đơn $\{X\}$ chỉ có phần tử là $X$. Do vậy, $\{X\} \subseteq \{X; Y\}$, dẫn đến $\{X\} \in \mathfrak{P}(\{X; Y\})$. Tương tự, $\{Y\} \in \mathfrak{P}(\{X; Y\})$.

Hiển nhiên $\{X; Y\} = \{X; Y\}$. Do vậy, $\{X; Y\} \in \mathfrak{P}(\{X; Y\})$.

Từ ba kết luận trên, chúng ta có $\mathfrak{P}(\{X; Y\}) \supseteq \{\varnothing; \{X\}; \{Y\}; \{X; Y\}\}$.

Mặt khác, với mọi $y \in \mathfrak{P}(\{X; Y\})$, theo định nghĩa có $y \subseteq \{X; Y\}$, hiểu theo một cách khác, $y$ chỉ nhận $X$ và $Y$ làm phần tử. Do mỗi phần tử $X$ và $Y$ chỉ có thể thuộc hoặc không thuộc $y$, nên các khả năng có thể xảy ra bao gồm\footnote{Tác giả sử dụng mở rộng của quy tắc chia trường hợp từ giải tích mệnh đề: ``Cho $P$ và $Q$ là hai mệnh đề, khi này, chúng ta có $P \land Q \lor P \land \neg Q \lor \neg P \land Q \lor \neg P \land \neg Q$.''.}:

\begin{itemize}
    \item \textcolor{colorEmphasis}{$X \in y$ và $Y \in y$}. Điều này cho chúng ta $\{X; Y\} \subseteq y$. Suy ra, $y = \{X; Y\}$.
    \item \textcolor{colorEmphasisCyan}{$X \in y$ và $Y \notin y$}. Gọi $w$ là một phần tử thuộc $y$. Có $\left[\begin{array}{l} w = X \\ w = Y \end{array}\right.$. Nhưng do $Y \notin y$ nên $w \neq Y$. Vậy $w = X$. Tóm lại, chúng ta có
          $$\forall w, \left(w \in y \iff w = X\right)$$
          và điều này cho $y = \{X\}$.
    \item \textcolor{colorEmphasisGreen}{$X \notin y$ và $Y \in y$}. Tương tự, có thể kết luận $y = \{Y\}$.
    \item \textcolor{colorEmphasisPurple}{$X \notin y$ và $Y \notin y$}. Giả sử tồn tại $w \in y$. Khi đó, cần phải có $w = X \lor w = Y$. Nhưng do $X \notin y$ và $Y \notin y$, nên điều này là vô lí. Vậy, $\forall w, w \notin y$. Suy ra, $y = \varnothing$ theo định nghĩa tập rỗng.
\end{itemize}

Qua bốn trường hợp, có $y \in \{\varnothing; \{X\}; \{Y\}; \{X; Y\}\}$. Do đó, $\mathfrak{P}(\{X; Y\}) \subseteq \{\varnothing; \{X\}; \{Y\}; \{X; Y\}\}$.

Kết hợp hai kết quả bao hàm, suy ra $\mathfrak{P}(\{X; Y\}) = \{\varnothing; \{X\}; \{Y\}; \{X; Y\}\}$. Điều phải chứng minh.

\begin{figure}[H]
    \begin{center}
        \begin{tikzpicture}[>=stealth]
            \draw[thick] (-3,0) ellipse (2.1cm and 3cm);
            \node at (-3, 3.3) {Tập hợp $\{a; b\}$};

            \node[circle, fill, inner sep=1.2pt, label=left:$a$] (a_left) at (-4, 0) {};
            \node[circle, fill, inner sep=1.2pt, label=left:$b$] (b_left) at (-2.5, -1.2) {};

            \draw[thick, dashed, colorEmphasis] (-2.5, 1.5) circle (0.4cm);
            \node[colorEmphasis] at (-2.5, 1.5) {$\varnothing$};

            \draw[thick, dashed, colorEmphasisCyan] (-4, 0) ellipse (0.5cm and 0.4cm);
            \node[colorEmphasisCyan] at (-4, 0.7) {$\{a\}$};

            \draw[thick, dashed, colorEmphasisGreen] (-2.5, -1.2) ellipse (0.4cm and 0.5cm);
            \node[colorEmphasisGreen] at (-1.8, -1.2) {$\{b\}$};

            \draw[thick, dashed, colorEmphasisPurple, rotate around={50:(-3.15,-0.45)}] (-3.15, -0.45) ellipse (1cm and 2cm);
            \node[colorEmphasisPurple] at (-4.2, -1.45) {$\{a; b\}$};

            \draw[->, ultra thick] (-0.5, 0) -- (1.9, 0) node[midway, above] {$\mathfrak{P}\left(\{a; b\}\right)$};

            \draw[thick] (5.5,0) ellipse (3.2cm and 3.5cm);
            \node at (5.5, 3.9) {Tập lũy thừa $\mathfrak{P}\left(\{a; b\}\right)$};

            % Tập rỗng
            \draw[thick, colorEmphasis] (5.5, 2) circle (0.7cm);
            \node[colorEmphasis] at (5.5, 2) {$\varnothing$};
            \draw[->, thick, colorEmphasis, dashed, shorten >= 5pt, shorten <= 5pt] (-2.2, 1.7) to[bend left=15] (4.85, 2);

            % Tập {a}
            \draw[thick, colorEmphasisCyan] (3.6, 0) ellipse (0.8cm and 0.6cm);
            \node[circle, fill, inner sep=1.2pt, label=below:$a$] (a_right1) at (3.6, 0) {};
            \node[colorEmphasisCyan] at (3.6, 0.9) {$\{a\}$};
            \draw[->, thick, colorEmphasisCyan, dashed, shorten >= 5pt, shorten <= 5pt] (-3.6, 0) to[bend right=25] (3.2, -0.5);

            % Tập {b}
            \draw[thick, colorEmphasisGreen] (7.4, 0) ellipse (0.6cm and 0.8cm);
            \node[circle, fill, inner sep=1.2pt, label=below:$b$] (b_right1) at (7.4, 0) {};
            \node[colorEmphasisGreen] at (8.3, 0) {$\{b\}$};
            \draw[->, thick, colorEmphasisGreen, dashed, shorten >= 8pt, shorten <= 8pt] (-2.4, -1.5) to[out=-60, in=180] (7, -0.1);

            % Tập {a, b}
            \draw[thick, colorEmphasisPurple, rotate around={-40:(5.5, -2)}] (5.5, -2) ellipse (1.2cm and 0.8cm);
            \node[circle, fill, inner sep=1.2pt, label=left:$a$] (a_right2) at (5.1, -1.7) {};
            \node[circle, fill, inner sep=1.2pt, label=right:$b$] (b_right2) at (5.9, -2.3) {};
            \node[colorEmphasisPurple] at (6.5, -1.2) {$\{a; b\}$};
            \draw[->, thick, colorEmphasisPurple, dashed, shorten >= 8pt, shorten <= 8pt] (-1.7, -1) to[out=50, in=-150] (5.1, -2.6);

        \end{tikzpicture}
    \end{center}
    \caption{Minh họa tiên đề tập lũy thừa với $\mathfrak{P}\left(\{a; b\}\right) = \big\{\varnothing; \{a\}; \{b\}; \{a; b\}\big\}$.}
\end{figure}


\stepcounter{subexercise}
\arabic{subexercise}. Giả sử $X \subseteq Y$. Gọi $w \in \mathfrak{P}(X)$. Theo định nghĩa, $w \subseteq X$. Theo tính chất bắc cầu của quan hệ bao hàm, có $w \subseteq Y$. Do đó, $w \in \mathfrak{P}(Y)$.

Từ đó, có $X \subseteq Y \implies \mathfrak{P}(X) \subseteq \mathfrak{P}(Y)$. Điều phải chứng minh.

\stepcounter{subexercise}
\arabic{subexercise}. Gọi $w$ là một phần tử bất kì thuộc $\mathfrak{P}(X) \cup \mathfrak{P}(Y)$. Theo định nghĩa của phép hợp, có $w \in \mathfrak{P}(X)$ hoặc $w \in \mathfrak{P}(Y)$.
Nếu $w \in \mathfrak{P}(X)$, theo định nghĩa tập lũy thừa, $w \subseteq X$. Mặt khác, $X \subseteq X \cup Y$, nên theo tính chất bắc cầu, $w \subseteq X \cup Y$, dẫn đến $w \in \mathfrak{P}(X \cup Y)$. Chứng minh hoàn toàn tương tự cho trường hợp $w \in \mathfrak{P}(Y)$, chúng ta cũng thu được $w \in \mathfrak{P}(X \cup Y)$.

Vậy, nhìn chung, chúng ta có
$$\forall w, \left( w \in \mathfrak{P}(X) \cup \mathfrak{P}(Y) \implies w \in \mathfrak{P}(X \cup Y) \right)$$
và do đó $\mathfrak{P}(X \cup Y) \supseteq \mathfrak{P}(X) \cup \mathfrak{P}(Y)$. Điều phải chứng minh.

